%%%%%%%%%%%%%%%%
%%% lev 21 %%%
%%%%%%%%%%%%%%%%

\Chapter{21}

\Verse{1}
{
	\hP{וַיֹּאמֶר יְהֹוָה אֶל מֹשֶׁה אֱמֹר אֶל הַכֹּהֲנִים בְּנֵי אַהֲרֹן וְאָמַרְתָּ אֲלֵהֶם לְנֶפֶשׁ לֹא יִטַּמָּא בְּעַמָּיו׃}
}
{
	\eP{
		And YHWH said to Moses, “Say to the priests, the sons of
		Aaron, and say to them: One shall not become impure for a
		person among his people
	}
}
{
%	\footnote{
%		become impure for a person. Meaning: to become impure by coming into
%		contact with a dead person.
%	}
}

\Verse{2}
{
	\hP{כִּי אִם לִשְׁאֵרוֹ הַקָּרֹב אֵלָיו לְאִמּוֹ וּלְאָבִיו וְלִבְנוֹ וּלְבִתּוֹ וּלְאָחִיו׃}
}
{
	\eP{
		except for his relative who is close to him: for his
		mother and for his father and for his son and for his
		daughter and for his brother
	}
}
{
}

\Verse{3}
{
	\hP{וְלַאֲחֹתוֹ הַבְּתוּלָה הַקְּרוֹבָה אֵלָיו אֲשֶׁר לֹא הָיְתָה לְאִישׁ לָהּ יִטַּמָּא׃}
}
{
	\eP{
		and for his virgin sister who is close to him, who has not
		been a husband’s; he may become impure for her.
	}
}
{
}

\Verse{4}
{
	\hP{לֹא יִטַּמָּא בַּעַל בְּעַמָּיו לְהֵחַלּוֹ׃}
}
{
	\eP{
		He shall not become impure as a husband among his people,
		for him to be desecrated.
	}
}
{
}

\Verse{5}
{
	\hP{לֹא [יִקְרְחוּ] (יקרחה) קרְחָה בְּרֹאשָׁם וּפְאַת זְקָנָם לֹא יְגַלֵּחוּ וּבִבְשָׂרָם לֹא יִשְׂרְטוּ שָׂרָטֶת׃}
}
{
	\eP{
		They shall not make a bald place on their head and shall
		not shave the corner of their beard and not make a cut in
		their flesh.
	}
}
{
	Qere Ketiv.
}

\Verse{6}
{
	\hP{קְדֹשִׁים יִהְיוּ לֵאלֹהֵיהֶם וְלֹא יְחַלְּלוּ שֵׁם אֱלֹהֵיהֶם כִּי אֶת אִשֵּׁי יְהֹוָה לֶחֶם אֱלֹהֵיהֶם הֵם מַקְרִיבִם וְהָיוּ קֹדֶשׁ׃}
}
{
	\eP{
		They shall be holy to their God and shall not desecrate
		their God’s name, because they are bringing forward YHWH’s
		offerings by fire, their God’s bread, so they shall be a
		holy thing.
	}
}
{
%	\footnote{
%		bread. Hebrew leem here refers broadly to the food of the sacrifices.
%	}
}

\Verse{7}
{
	\hP{אִשָּׁה זֹנָה וַחֲלָלָה לֹא יִקָּחוּ וְאִשָּׁה גְּרוּשָׁה מֵאִישָׁהּ לֹא יִקָּחוּ כִּי קָדֹשׁ הוּא לֵאלֹהָיו׃}
}
{
	\eP{
		They shall not take a woman who is a prostitute and
		desecrated, and they shall not take a woman who is
		divorced from her husband, because he is holy to his God.
	}
}
{
}

\Verse{8}
{
	\hP{וְקִדַּשְׁתּוֹ כִּי אֶת לֶחֶם אֱלֹהֶיךָ הוּא מַקְרִיב קָדֹשׁ יִהְיֶה לָּךְ כִּי קָדוֹשׁ אֲנִי יְהֹוָה מְקַדִּשְׁכֶם׃}
}
{
	\eP{
		And you shall regard him as holy because he is bringing
		forward your God’s bread. He shall be holy to you because
		I, YHWH, who makes you holy, am holy.
	}
}
{
%	\footnote{
%		holy … holy … holy. The word occurs seven times in three verses, and
%		four times just in this last verse. It is not redundant, but rather
%		emphatic. This entire section of the Torah is concerned with holiness—it
%		is referred to in biblical scholarship as the Holiness Code—and in these
%		verses it is focused on the persons who are designated to ensure and
%		carry out that holiness and to be holy themselves. Footnote 21:8. He
%		shall be holy … I make you holy. Note that there are degrees of
%		holiness. All the people are holy, but the priests still are holy in a
%		singular way because of their role. This matter of degree between the
%		people’s holiness and the priests’ holiness will become an issue in the
%		matter of Korah’s rebellion (Num 16:3–7).
%	}
}

\Verse{9}
{
	\hP{וּבַת אִישׁ כֹּהֵן כִּי תֵחֵל לִזְנוֹת אֶת אָבִיהָ הִיא מְחַלֶּלֶת בָּאֵשׁ תִּשָּׂרֵף׃}
}
{
	\eP{
		And the daughter of a man who is a priest who will become
		desecrated, to whore: she is desecrating her father. She
		shall be burned in fire.
	}
}
{
}

\Verse{10}
{
	\hP{וְהַכֹּהֵן הַגָּדוֹל מֵאֶחָיו אֲשֶׁר יוּצַק עַל רֹאשׁוֹ שֶׁמֶן הַמִּשְׁחָה וּמִלֵּא אֶת יָדוֹ לִלְבֹּשׁ אֶת הַבְּגָדִים אֶת רֹאשׁוֹ לֹא יִפְרָע וּבְגָדָיו לֹא יִפְרֹם׃}
}
{
	\eP{
		“And the priest who is the most senior of his brothers, on
		whose head the anointing oil will be poured and who has
		filled his hand to wear the clothes, he shall not let
		loose the hair of his head and not tear his clothes
	}
}
{
%	\footnote{
%		most senior. Literally, “the biggest priest,” Hebrew hakkhn haggdôl. It
%		is the standard term for the high priest and could therefore also be
%		translated “the high priest from among his brothers.” Israel’s first two
%		high priests, Aaron and Eliezer, are each the oldest living son; but the
%		high priests in the following centuries may or may not be the eldest
%		sons of the preceding high priests. Footnote 21:10. filled his hand. See
%		comment on Exod 29:29.
%	}
}

\Verse{11}
{
	\hP{וְעַל כּל נַפְשֹׁת מֵת לֹא יָבֹא לְאָבִיו וּלְאִמּוֹ לֹא יִטַּמָּא׃}
}
{
	\eP{
		and he shall not come to any dead persons; he shall not
		become impure for his father and for his mother.
	}
}
{
}

\Verse{12}
{
	\hP{וּמִן הַמִּקְדָּשׁ לֹא יֵצֵא וְלֹא יְחַלֵּל אֵת מִקְדַּשׁ אֱלֹהָיו כִּי נֵזֶר שֶׁמֶן מִשְׁחַת אֱלֹהָיו עָלָיו אֲנִי יְהֹוָה׃}
}
{
	\eP{
		And he shall not go out from the holy place, so he will
		not desecrate his God’s holy place, because the crown with
		his God’s anointing oil is on him. I am YHWH.
	}
}
{
%	\footnote{
%		crown with anointing oil. It is unclear whether this term (Hebrew nzer)
%		refers to the priest’s crown or to his hair (the crown of his head). The
%		high priest’s crown is associated with the anointing in Exod 29:6–7; but
%		elsewhere in Priestly law this term refers to consecrated hair (Num
%		6:19). Here it is consecrated by the anointing oil. There it is
%		consecrated by a vow taken by a Nazirite.
%	}
}

\Verse{13}
{
	\hP{וְהוּא אִשָּׁה בִבְתוּלֶיהָ יִקָּח׃}
}
{
	\eP{“And he shall take a woman in her virginity.}
}
{
}

\Verse{14}
{
	\hP{אַלְמָנָה וּגְרוּשָׁה וַחֲלָלָה זֹנָה אֶת אֵלֶּה לֹא יִקָּח כִּי אִם בְּתוּלָה מֵעַמָּיו יִקַּח אִשָּׁה׃}
}
{
	\eP{
		A widow and a divorcée and a desecrated woman, a
		prostitute: he shall not take these; he shall take none
		except a virgin from his people.
	}
}
{
}

\Verse{15}
{
	\hP{וְלֹא יְחַלֵּל זַרְעוֹ בְּעַמָּיו כִּי אֲנִי יְהֹוָה מְקַדְּשׁוֹ׃}
}
{
	\eP{
		And he shall not desecrate his seed among his people,
		because I, YHWH, make him holy.”
	}
}
{
}

\Verse{16}
{
	\hP{וַיְדַבֵּר יְהֹוָה אֶל מֹשֶׁה לֵּאמֹר׃}
}
{
	\eP{And YHWH spoke to Moses, saying,}
}
{
}

\Verse{17}
{
	\hP{דַּבֵּר אֶל אַהֲרֹן לֵאמֹר אִישׁ מִזַּרְעֲךָ לְדֹרֹתָם אֲשֶׁר יִהְיֶה בוֹ מוּם לֹא יִקְרַב לְהַקְרִיב לֶחֶם אֱלֹהָיו׃}
}
{
	\eP{
		“Speak to Aaron, saying: A man from your seed, through
		their generations, in whom will be an injury shall not
		come forward to bring forward his God’s bread.
	}
}
{
%	\footnote{
%		injury. Meaning: a hurt or impairment to his body. This has commonly
%		been understood as “blemish” or “defect,” but those terms may connote
%		something like a shortcoming in the person, denigrating the person as
%		being marred or unworthy. But the actual list indicates that these are
%		injuries or impairments that occurred either at birth or later in the
%		person’s life. They do not reflect on the character or quality of the
%		person. But they do disqualify the person from offering sacrifices. For
%		examples of what constitutes an injury, see 21:18–20 and 24:19–20.
%	}
}

\Verse{18}
{
	\hP{כִּי כל אִישׁ אֲשֶׁר בּוֹ מוּם לֹא יִקְרָב אִישׁ עִוֵּר אוֹ פִסֵּחַ אוֹ חָרֻם אוֹ שָׂרוּעַ׃}
}
{
	\eP{
		Because any man in whom is an injury shall not come
		forward: a blind man or a cripple or mutilated or who has
		an elongated limb
	}
}
{
}

\Verse{19}
{
	\hP{אוֹ אִישׁ אֲשֶׁר יִהְיֶה בוֹ שֶׁבֶר רָגֶל אוֹ שֶׁבֶר יָד׃}
}
{
	\eP{
		or a man who has a break of the leg or a break of the arm
		in him
	}
}
{
}

\Verse{20}
{
	\hP{אוֹ גִבֵּן אוֹ דַק אוֹ תְּבַלֻּל בְּעֵינוֹ אוֹ גָרָב אוֹ יַלֶּפֶת אוֹ מְרוֹחַ אָשֶׁךְ׃}
}
{
	\eP{
		or a hunchback or a dwarf or who has a spotting in his eye
		or is scabbed or scurvied or has crushed testicles.
	}
}
{
}

\Verse{21}
{
	\hP{כּל אִישׁ אֲשֶׁר בּוֹ מוּם מִזֶּרַע אַהֲרֹן הַכֹּהֵן לֹא יִגַּשׁ לְהַקְרִיב אֶת אִשֵּׁי יְהֹוָה מוּם בּוֹ אֵת לֶחֶם אֱלֹהָיו לֹא יִגַּשׁ לְהַקְרִיב׃}
}
{
	\eP{
		Any man from the seed of Aaron the priest in whom is an
		injury shall not come near to bring forward YHWH’s
		offerings by fire. He has an injury in him. He shall not
		come near to bring forward his God’s bread.
	}
}
{
}

\Verse{22}
{
	\hP{לֶחֶם אֱלֹהָיו מִקּדְשֵׁי הַקֳּדָשִׁים וּמִן הַקֳּדָשִׁים יֹאכֵל׃}
}
{
	\eP{
		He shall eat his God’s bread, from the holies of holies
		and from the holies.
	}
}
{
}

\Verse{23}
{
	\hP{אַךְ אֶל הַפָּרֹכֶת לֹא יָבֹא וְאֶל הַמִּזְבֵּחַ לֹא יִגַּשׁ כִּי מוּם בּוֹ וְלֹא יְחַלֵּל אֶת מִקְדָּשַׁי כִּי אֲנִי יְהֹוָה מְקַדְּשָׁם׃}
}
{
	\eP{
		Just: he shall not come to the pavilion, and he shall not
		come near to the altar, because he has an injury in him.
		And he shall not desecrate my holy places, because I,
		YHWH, make them holy.”
	}
}
{
%	\footnote{
%		he shall not come near to the altar, because he has an injury in him.
%		And he shall not desecrate my holy places. What should we make of the
%		fact that someone who is injured physically is disqualified from various
%		priestly services? As I indicated in the preceding comment, it is not a
%		denigration of the person’s worth or of his ethical merit. It is rather
%		another case that relates to the ritual versus the ethical realm. At the
%		time that I am writing this, we have lost the sense of zones of
%		holiness, impurity, and so on. The Torah contains the element of awe for
%		the sacred, inspiring great care regarding who enters the zone, what
%		fits the zone, and what the zone’s boundaries are. By present standards,
%		most would judge it to be unkind and hurtful to exclude an injured
%		person from the priesthood. The same standard applies in the biblical
%		world in ethical matters: one must be considerate of persons who are
%		handicapped (Lev 19:14). But in the ritual matter of the priesthood,
%		physical conditions supersede such considerations.
%	}
}

\Verse{24}
{
	\hP{וַיְדַבֵּר מֹשֶׁה אֶל אַהֲרֹן וְאֶל בָּנָיו וְאֶל כּל בְּנֵי יִשְׂרָאֵל׃}
}
{
	\eP{
		So Moses spoke to Aaron and to his sons and to all the
		children of Israel.
	}
}
{
}
